\clearpage\section*{从想法到可以使用的结果}\label{guide}\markboth{工作框架}{}
\lede{用自然语言说明目标，让编程智能体读取项目、修改文件、运行检查，再根据实际结果继续完善。智能体会反复调用工具，根据结果决定下一步，朝目标推进。}
\diagram{p_route}
\topic{一条主线，按需要增加方法}
先说清目标，调研已有方案，再让 AI 实现、检查和修正。把需求、依据与进度留在项目文件中，按需要增加协作与自动化。重点是理解每种方法解决什么问题、何时使用；具体工具操作可以让 AI 查证后协助完成。

\begin{tabularx}{\linewidth}{@{}L{35mm}Y@{}}
\toprule 当前情况 & 对应做法\\\midrule
想法还比较模糊 & 澄清目标、范围和完成标准，先形成最小可用版本的计划\\
开始开发或修改 & 读取项目与相关资料，复用已有做法，小范围实现并及时运行\\
结果不符合预期 & 按真实操作验收，保留出错记录，先定位原因再修复\\
有多项独立工作 & 可同时做只读调查；统一安排改动，重要成果再独立审查\\
任务中断或重复出现 & 保存进度；把有效方法写成规则或技能，再考虑批量执行\\
\bottomrule\end{tabularx}
\topic{每章浅色框里的提示词怎么用}
\begin{enumerate}
\item \textbf{找到并复制。}每章正文中的浅色大框就是提示词。按任务选一框，复制标题下的全部文字。
\item \textbf{替换方括号。}把“[目标]”等内容换成自己的要求，例如“添加待办事项，刷新后仍能看到”。
\item \textbf{粘贴并发送。}在 Codex、Claude Code 中打开项目，把改好的整段文字发到对话框，再按步骤运行并检查结果。
\end{enumerate}
\topic{人的判断放在哪里}
人决定目标、取舍和是否接受结果；AI 比较方案、执行并提供证据。这里的判断力（taste）指理解用户并确定优先级。演示版可侧重表达，日常工具可侧重可靠性；先满足必要要求，再权衡体验、成本与速度。用场景、样例与可观察标准表达偏好，并检验标准是否合适。
